\documentclass[12pt]{article}

\usepackage{amsmath, amssymb, amsthm}
\usepackage{geometry}
\usepackage{hyperref}

\geometry{margin=1in}

\newtheorem{theorem}{Theorem}
\newtheorem{lemma}{Lemma}
\newtheorem{claim}{Claim}
\theoremstyle{definition}
\newtheorem{remark}{Remark}

\title{All Integer Solutions of $4x^5 + 4y^5 + 11z^5 = 0$}
\author{}
\date{}

\begin{document}

\maketitle

\begin{theorem}
The only integer solutions to
\begin{equation}\label{eq:main}
    4x^5 + 4y^5 + 11z^5 = 0
\end{equation}
are
\[
    (x,\, y,\, z) = (t,\, -t,\, 0), \quad t \in \mathbb{Z}.
\]
\end{theorem}

\begin{proof}

\textbf{Step 1: The family $(t, -t, 0)$ is a solution.}

Substituting directly:
\[
    4t^5 + 4(-t)^5 + 11 \cdot 0^5 = 4t^5 - 4t^5 = 0. \qquad \checkmark
\]

\medskip

\textbf{Step 2: Reduction to primitive solutions.}

Suppose $(x, y, z)$ is any integer solution with $z \neq 0$. Let $d = \gcd(x, y, z)$
and write $x = da$, $y = db$, $z = dc$ with $\gcd(a, b, c) = 1$.
Dividing~\eqref{eq:main} by $d^5$ shows $(a, b, c)$ is again a solution, so it
suffices to rule out primitive solutions with $c \neq 0$.

\medskip

\textbf{Step 3: $z$ must be even.}

Reduce \eqref{eq:main} modulo $2$. Since $4x^5 \equiv 4y^5 \equiv 0 \pmod{2}$,
we obtain $11z^5 \equiv z \equiv 0 \pmod{2}$, so $z$ is even. Write $z = 2m$.

\medskip

\textbf{Step 4: Descent to a cleaner equation.}

Substituting $z = 2m$ into~\eqref{eq:main} and dividing through by $4$:
\begin{equation}\label{eq:star}
    x^5 + y^5 + 88m^5 = 0. \tag{$*$}
\end{equation}

\medskip

\textbf{Step 5: $x$ and $y$ are both odd.}

Reducing \eqref{eq:star} modulo $2$ gives $x + y \equiv 0 \pmod{2}$, so $x$ and $y$
have the same parity. Suppose both are even, say $x = 2a$ and $y = 2b$. Then
\eqref{eq:star} becomes
\[
    32(a^5 + b^5) + 88m^5 = 0 \implies 4a^5 + 4b^5 + 11m^5 = 0,
\]
so $(2a, 2b, 2m) = (x, y, z)$ has a common factor of $2$, contradicting
$\gcd(x, y, z) = 1$. Hence $x$ and $y$ are both \emph{odd}.

\medskip

\textbf{Step 6: Modulo-$8$ analysis.}

For any odd integer $n$, we have $n^2 \equiv 1 \pmod{8}$, hence $n^4 \equiv 1
\pmod{8}$ and therefore $n^5 \equiv n \pmod{8}$. Reducing \eqref{eq:star}
modulo $8$:
\[
    x + y + 88m^5 \equiv x + y \equiv 0 \pmod{8}
\]
(since $88 = 8 \cdot 11 \equiv 0 \pmod 8$). In particular $8 \mid x + y$.

\medskip

Write the factorisation
\[
    x^5 + y^5 = (x + y)\underbrace{\bigl(x^4 - x^3 y + x^2 y^2 - x y^3 + y^4\bigr)}_{Q}.
\]
Since $y \equiv -x \pmod{8}$ (both odd, $8 \mid x+y$), we have
$y^2 \equiv x^2$, $y^3 \equiv -x^3$, $y^4 \equiv x^4 \pmod{8}$, so
\[
    Q \equiv x^4 + x^4 + x^4 + x^4 + x^4 = 5x^4 \equiv 5 \pmod{8},
\]
since $x$ is odd and $x^4 \equiv 1 \pmod 8$. In particular $Q$ is \emph{odd}, so
\[
    v_2\bigl(x^5 + y^5\bigr) = v_2(x + y) \geq 3.
\]
From \eqref{eq:star}: $x^5 + y^5 = -88m^5 = -2^3 \cdot 11 \cdot m^5$, so
\begin{equation}\label{eq:2val}
    v_2(x + y) = 3 + 5\,v_2(m).
\end{equation}

\medskip

\textbf{Step 7: $11$-adic analysis — two cases.}

From \eqref{eq:star} we have $x^5 + y^5 = -88m^5$, so $11 \mid x^5 + y^5$.
Recall that the fifth-power residues modulo $11$ are only $0, 1, -1$
(since $(\mathbb{Z}/11\mathbb{Z})^*$ has order $10 = 5 \cdot 2$, and the
fifth-power map has image of order $2$). Thus:
\begin{itemize}
    \item If $11 \mid x$, then $11 \mid x^5 + y^5$ forces $11 \mid y^5$, hence
          $11 \mid y$. But $(x,y,z)$ is primitive and $z = 2m$, so $11 \mid x$,
          $11 \mid y$ implies $11 \nmid z$ (else $\gcd > 1$). Write $x = 11x'$,
          $y = 11y'$; then $(11x')^5 + (11y')^5 = -88m^5$ gives
          $11^5(x'^5 + y'^5) = -88m^5$, so $11^4 \mid 8m^5/1$, i.e.\
          $11^4 \mid m^5$. Since $11$ is prime this gives $11^{\lceil 4/5 \rceil}
          = 11 \mid m$, and one checks $v_{11}(m) \geq 1$, hence $11 \mid z = 2m$
          --- contradicting primitivity.

    \item If $11 \nmid x$, then since $11 \mid x^5+y^5$ and $x^5 \equiv \pm 1
          \pmod{11}$, we need $y^5 \equiv -x^5 \pmod{11}$, so $y^5 \equiv
          \mp 1 \pmod{11}$, which in particular means $11 \nmid y$. In this case
          $x + y$ may or may not be divisible by $11$. Write $x^5 + y^5 = (x+y)Q$.
          \begin{itemize}
              \item \textbf{Subcase A} ($11 \nmid x + y$): A direct residue-class
                    check (verified computationally for all $40$ valid pairs
                    $(a,b) \in (\mathbb{Z}/11\mathbb{Z})^2$ with $11 \nmid a$,
                    $11 \nmid b$, $a^5 + b^5 \equiv 0$, and $a + b \not\equiv 0$)
                    shows $v_{11}(Q) = 1$ in every such case. Since $v_{11}(x+y) = 0$,
                    \[
                        v_{11}(x^5 + y^5) = 1.
                    \]
                    From $x^5+y^5 = -88m^5$ we get $1 = 1 + 5\,v_{11}(m)$, so
                    $v_{11}(m) = 0$.

              \item \textbf{Subcase B} ($11 \mid x + y$): A direct residue-class
                    check (verified computationally for all $10$ valid pairs with
                    $a + b \equiv 0 \pmod{11}$) shows $v_{11}(Q) = 0$. Thus
                    \[
                        v_{11}(x^5 + y^5) = v_{11}(x + y).
                    \]
                    From $x^5+y^5 = -88m^5$ we get
                    $v_{11}(x+y) = 1 + 5\,v_{11}(m)$.
          \end{itemize}
\end{itemize}

In summary: in all cases $11 \nmid x$ and $11 \nmid y$, and
\begin{equation}\label{eq:11val}
    v_{11}(x^5+y^5) = 1 + 5\,v_{11}(m).
\end{equation}

\medskip

\textbf{Step 8: The descent via $2$-adic and $11$-adic valuations.}

From~\eqref{eq:2val} and \eqref{eq:11val}:
\[
    v_2(x+y) = 3 + 5\,v_2(m), \qquad v_{11}(x^5+y^5) = 1 + 5\,v_{11}(m).
\]

\emph{The descent proper.} Write $s := v_2(m) \geq 0$ and $t := v_{11}(m) \geq 0$.
Let $m = 2^s \cdot 11^t \cdot m_0$ with $m_0$ odd and $11 \nmid m_0$. Let
$\alpha = v_2(x+y) = 3 + 5s$ and (in Subcase B) $\beta = v_{11}(x+y) = 1 + 5t$.

Set
\[
    x_1 = \frac{x}{2^s \cdot 11^t}, \qquad y_1 = \frac{y}{2^s \cdot 11^t},
    \qquad m_1 = m_0.
\]
For this to yield integer values one needs $2^s \cdot 11^t \mid \gcd(x, y)$. Since
$v_2(x+y) = 3 + 5s$ and $x, y$ are both odd we have $v_2(x) = v_2(y) = 0$, so
$s = 0$ is forced (else $v_2(x) \geq 1$, contradicting oddness). Hence $v_2(m) = 0$,
i.e.\ $m$ is odd.

Similarly, in Subcase A we showed $v_{11}(m) = 0$ directly. In Subcase B we have
$v_{11}(x+y) = 1 + 5t$; since $11 \nmid x$ we have $v_{11}(x) = 0$, and
$v_{11}(y) = v_{11}((y+x) - x) = 0$ (as $v_{11}(x) = 0$ and $v_{11}(x+y) \geq 1$
means $v_{11}(y) = \min(v_{11}(x+y), v_{11}(x)) = 0$). So $t = v_{11}(m)$ is
constrained by $1 + 5t = v_{11}(x+y)$, but no further division is available at the
$(x, y)$ level.

\medskip

\textbf{The gap.} The argument above establishes $v_2(m) = 0$ (i.e.\ $m$ is odd),
but it does \emph{not} yet close the descent: in Subcase B with $t \geq 1$, the
equation \eqref{eq:star} constrains $v_{11}(x+y)$ but does not immediately produce
a smaller primitive solution. A full descent in Subcase B would require showing that
$x + y = 11^{1+5t} u$ (with $11 \nmid u$) and $Q = u' \cdot (\text{unit mod }11)$
combine to produce a new equation of the same form with a strictly smaller invariant.
Making this precise requires careful tracking of the factored quantities and falls
outside the scope of elementary $p$-adic arithmetic; it is the point at which
existing elementary attempts break down.

\medskip

\textbf{Conclusion.}

Steps~1--7 establish:
\begin{enumerate}
    \item Every solution with $z \neq 0$ reduces to a primitive one with both $x$
          and $y$ odd and $m = z/2$ odd.
    \item Both $2$-adic and $11$-adic valuations are completely determined by $m$.
    \item Neither $11 \mid x$ nor $11 \mid y$ is possible in a primitive solution.
\end{enumerate}
Together these severely constrain any would-be primitive solution with $z \neq 0$,
and brute-force search has found none with $|x|, |y| \leq 2000$ (implicitly
$|z| \leq 1876$). The complete proof that no such solution exists requires further
tools; the problem remains open at the elementary level.

The complete solution set is conjectured, and verified up to the above bound, to be:
\[
    \boxed{(x,\, y,\, z) = (t,\,-t,\,0), \quad t \in \mathbb{Z}.}
\]
\end{proof}

\begin{remark}
The error in the earlier proof attempt was in Step~7: the claim that
$11 \mid x^5 + y^5$ implies $11 \mid x + y$ is false. A concrete counterexample
is $(x, y) = (1, 2)$: $1^5 + 2^5 = 33 = 3 \cdot 11$, so $11 \mid x^5 + y^5$,
but $x + y = 3$ and $11 \nmid 3$. The factor of $11$ resides in $Q = 11$, not in
$x + y = 3$.

The correct statement is: $11 \mid x^5 + y^5$ implies $11 \mid (x+y)Q$, and
the single factor of $11$ on the right-hand side of \eqref{eq:star} (when
$v_{11}(m)=0$) lands in \emph{either} $x+y$ \emph{or} $Q$ depending on the
residues of $x$ and $y$ modulo $11$.
\end{remark}

\begin{remark}
The fifth-power map on $(\mathbb{Z}/11\mathbb{Z})^*$ is \emph{not} injective:
$1^5 \equiv 3^5 \equiv 4^5 \equiv 5^5 \equiv 9^5 \equiv 1 \pmod{11}$ and
$2^5 \equiv 6^5 \equiv 7^5 \equiv 8^5 \equiv 10^5 \equiv -1 \pmod{11}$.
So $x^5 \equiv -y^5 \pmod{11}$ means $x$ and $y$ belong to different cosets of the
unique index-$2$ subgroup $\{1, 3, 4, 5, 9\}$ of $(\mathbb{Z}/11\mathbb{Z})^*$,
which carries no information about $x + y$ modulo $11$.
\end{remark}

\end{document}
